\begin{abstract}
\noindent
Este relatório responde, um a um, aos oito objetivos específicos do subprojeto, a
partir de um experimento de $29$ bases $\times$ $10$ folds estratificados $\times$
$5$ modos de geração de pool ($1450$ folds, $0$ erros), com pools de $M=100$
classificadores. Três dos cinco modos separam o efeito do classificador-base do
efeito do método de geração --- um confounder de versões anteriores deste relatório
--- e um deles reativa a votação de medidas de complexidade do PGDCS completo,
antes desligada. As medidas \On{1} a \On{M} foram implementadas e verificadas; a
monotonicidade $\On{1}\ge\On{2}\ge\dots\ge\On{M}$ vale em $1450/1450$ folds, assim
como $\On{1}\ge$ votação majoritária. Os limites foram confrontados com a acurácia
individual, a votação majoritária, a média de probabilidades e \emph{doze} métodos
de seleção dinâmica de classificadores e de \emph{ensembles}.

Os quatro resultados centrais são: \textbf{(i)} o eixo que organiza os resultados é
o classificador-base, não o método de geração --- votar as medidas de complexidade
do PGDCS ou buscar com o algoritmo genético não muda \On{1}, $\mvr$ ou $\dfe$ de
forma estatisticamente sustentada ($p\ge0.13$ em todos os contrastes), e o padrão
inteiro (\On{1} alto, $\mvr$ baixo, folga grande) reaparece igual sem busca alguma,
some apenas quando o classificador-base muda para árvore; \textbf{(ii)} nenhum
nível intermediário de \On{N} serve como limite superior realista \emph{nos
problemas binários} --- o menor $N$ em que \On{N} cai abaixo da votação majoritária
fica em $N\approx M/2$ nos cinco modos (mediana $52$--$53$, Friedman $p=0.48$),
porque aí vale $\On{M/2+1}\le\mvr\le\On{M/2}$, verificado em $1100$ de $1100$ folds
binários; \textbf{(iii)} a folga $\On{1}-\mvr$ é previsível \emph{antes} de treinar
qualquer método de seleção, pelo índice de redundância de erros $\dfe$ (Spearman
$-0.89$, $n=145$), validado fora da amostra ($85{,}5\%$ de acerto
\emph{leave-one-dataset-out} contra $75{,}9\%$ da regra trivial); \textbf{(iv)}
mesmo tomando o melhor de dez métodos dinâmicos escolhido \emph{a posteriori} no
teste, a seleção dinâmica recupera entre $14{,}2\%$ e $24{,}3\%$ dessa folga,
conforme o modo --- o \On{1} é um limite inatingível na prática, não uma meta. O
relatório fecha com oito recomendações operacionais.
\end{abstract}
